\subsection{Final construction and decoder}
\label{sec:final}

All ingredients are now in place.  The final recursion uses the known-powers
gadgets $Q_{2^t-1}$, the fill gadget $A_{2^l}$, the $T_{k,2^l}$ remainder
recursion, and the three causal closure operations.  Each recursive branch
comes with its own cutoff-respecting decoder: the closure lemmas handle local
combinations, and the stage tables handle $T$.  The decoder summaries used
below are
\Cref{alg:decode-Q-2kminus1} for $Q_{2^t-1}$,
\Cref{alg:decode-fill} for $A_{2^l}$,
and \Cref{alg:decode-Rk2l} for the combined remainder polynomial used in the $T_{k,2^l}$ recursion;
they assemble into \Cref{alg:final-decoder}.  The schedules certified in
Lean substitute the peeled gadget $Q^{\mathrm{peel}}$ (\Cref{sec:peeled-Q})
for every known-powers instance; \Cref{thm:construction-height-peeled}
records the resulting $O(\log n)$ height at the identical gate ledger.
\Cref{fig:final-recursion} charts the routing of the strong induction.

\begin{figure}[H]
  \centering
  \resizebox{\textwidth}{!}{% Routing diagram for the strong induction and the final multiplication.
\begin{tikzpicture}[x=1cm,y=1cm,>=Latex]
  \node[fp active, rounded corners, minimum width=30mm, minimum height=8mm]
    (n) at (6.75,5.2) {odd target degree $n\ge3$};

  \node[fp known, rounded corners, minimum width=19mm] (seven) at (0.9,3.75)
    {$n=7$\\direct septic};
  \node[fp second, rounded corners, minimum width=30mm] (bases) at (3.45,3.75)
    {$3$ base; $15$ fused $k=1$\\$27,31$ septic bridges};
  \node[fp active, rounded corners, minimum width=23mm] (one) at (6.35,3.75)
    {$n\equiv1\pmod4$\\$T_{2k,2}$ crown};
  \node[fp active, rounded corners, minimum width=23mm] (three) at (9.05,3.75)
    {$n\equiv3\pmod8$\\$8k+3$ shell};
  \node[fp active, rounded corners, minimum width=23mm] (sevenmod) at (11.85,3.75)
    {$n\equiv7\pmod8$\\$8k+7$ shells};

  % Route through one branch bus: the five cases are mutually exclusive, and
  % the common trunk avoids the visually misleading crossings of five curves.
  \coordinate (route) at (6.75,4.55);
  \draw[line width=.65pt, draw=black!75] (n.south) -- (route)
    -- (0.9,4.55) -- (11.85,4.55);
  \draw[fp flow] (0.9,4.55) -- (seven.north);
  \draw[fp flow] (3.45,4.55) -- (bases.north);
  \draw[fp flow] (6.35,4.55) -- (one.north);
  \draw[fp flow] (9.05,4.55) -- (three.north);
  \draw[fp flow] (11.85,4.55) -- (sevenmod.north);

  \node[fp recursive, rounded corners, minimum width=28mm] (small3) at (8.75,2.1)
    {smaller pair at\\$m=(n+1)/4$};
  \node[fp recursive, rounded corners, minimum width=28mm] (small7) at (11.95,2.1)
    {smaller pair at\\$m=(n-3)/4$};
  \draw[fp flow] (small3.north) -- node[fp label, right] {insert} (three.south);
  \draw[fp flow] (small7.north) -- node[fp label, right] {insert} (sevenmod.south);
  \node[fp label, anchor=west] at (8.0,1.4)
    {$m<n$, and the explicit exceptions ensure $m\ne7$};

  \node[fp recursive, minimum width=73mm, minimum height=7mm] (pair) at (4.2,0.55)
    {shared construction of $(T_n^{(1)},T_n^{(2)})$: $(n-1)/2$ products};
  \node[fp gate] (mul) at (8.65,0.55) {$\times$};
  \node[fp active, minimum width=34mm, minimum height=7mm] (out) at (11.3,0.55)
    {$P_n=xT_n^{(1)}+T_n^{(2)}$\\$(n+1)/2$ products total};
  \node[fp label, below=1mm of mul] {multiply by $x$};
  \draw[fp flow] (pair.east) -- (mul.west);
  \draw[fp flow] (mul.east) -- (out.west);
\end{tikzpicture}
}
  \caption{Routing of the final strong induction.  The $4k+1$ branch closes
  directly through the $T$ recursion.  The other two congruence classes wrap
  a strictly smaller compatible pair in one or two square shells.  The
  exceptional list removes exactly the cost-optimal recursive calls that
  would otherwise land at the missing three-product degree-$7$ realization
  (and the degree-$15$ byproduct mismatch).
  The bottom row separates the cost of constructing the shared pair from the
  final multiplication by~$x$.  All known-powers instances inside these
  branches are scheduled as the peeled $Q^{\mathrm{peel}}$
  (\Cref{thm:construction-height-peeled}), giving height
  $2\lceil\log_2 n\rceil+4$ for the odd degrees charted here (the even
  lift of \Cref{thm:construction-count} adds one).}
  \label{fig:final-recursion}
\end{figure}

\begin{algorithm}[H]
  \caption{Final odd-degree construction}\label{alg:final-construction}
  \begin{algorithmic}
    \Require odd integer $n\ge 3$
    \Ensure a decodable monic polynomial $P_n$; if $n\ne7$, also a
    jointly $(n-1)/2$-realized splittable pair
    $(T^{(1)}_n,T^{(2)}_n)$ that is compatible with no auxiliary data and
    records the required power byproducts

    \If{$n=7$}
      \State Use the direct septic construction in \Cref{lem:septic-base}.
    \ElsIf{$n\in\{3,15,27,31\}$}
      \State Use the explicit special-case constructions in \Cref{sec:special-cases}.
    \ElsIf{$n\equiv 1\pmod 4$}
      \State Write $n=4k+1$ and construct $(T^{(1)}_{4k+1},T^{(2)}_{4k+1})$ via the $T_{2k,2}$ call; see \Cref{lem:4k+1-splittable}.
    \ElsIf{$n\equiv 3\pmod 8$}
      \State Write $n=8k+3$ and extend the smaller compatible joint realization for $2k+1$ by \Cref{alg:constr-8k+3}.
    \Else
      \State Write $n=8k+7$ with $k\ge 2$ and extend the smaller compatible joint realization for $2k+1$ by \Cref{alg:constr-8k+7} (using $\bar Q_{15}$ or $\bar Q_{8k+7}$ as needed).
    \EndIf
  \end{algorithmic}
\end{algorithm}

\begin{algorithm}[H]
  \caption{Final decoder for $P_n$ (structural recursion)}\label{alg:final-decoder}
  \begin{algorithmic}
    \Require odd $n\ge 3$ and $P_n$ produced by \Cref{alg:final-construction}.
      For the peeled schedules of \Cref{thm:construction-height-peeled}, every
      invocation of \Cref{alg:decode-Q-2kminus1} below is replaced by the
      decoder of \Cref{lem:peeled-Q-decodable}, which has the identical
      interface.
    \Ensure the full parameter list $(\alpha_0,\ldots,\alpha_{n-1})$

    \If{$n=7$}
      \State Apply the triangular decoder in \Cref{lem:septic-base}.
    \ElsIf{$n\in\{3,15,27,31\}$}
      \State Decode by the explicit special-case proofs in \Cref{sec:special-cases}.
    \ElsIf{$n\equiv 1\pmod 4$}
      \State Write $n=4k+1$.
      \State Use the five top-coefficient pivots in \Cref{lem:4k+1-splittable} to derive $H_2,H_4,\tilde H_4=H_4+\rho$ and the five outer parameters.
      \State Derive the pair $(U,V)=T_{k,4}(x,H_2,H_4,\tilde H_4)$ from $P_{4k+1}=xU+V$ using the cutoff decoder in \Cref{lem:4k+1-splittable}.
      \State Form $x(U-H_4^k)+(V-\tilde H_4^k)$ and apply \Cref{alg:decode-Rk2l} at $(k,l)=(k,2)$ to extract $\alpha_0,\ldots,\alpha_{4k-5}$.
    \ElsIf{$n\equiv 3\pmod 8$}
      \State Write $n=8k+3$ and apply the square-gadget extraction in \Cref{lem:8k+3-splittable} to recover the auxiliary polynomials and the squared smaller pair.
      \State Recover $P_{2k+1}$ by monic square roots and recurse on it; this reconstructs $H_2$.
      \State Decode $S^{(1)}_2=Q_{4k+1}$ first via \Cref{lem:Q4k+1-from-H2}, which reconstructs $H_4$; only then decode $S^{(1)}_3=\mathcal Q_{2k-1}$ by the dispatch of \Cref{lem:odd-gadgets-H2H4}, with \Cref{alg:decode-Q-2kminus1} for nested Mersenne blocks.
    \Else
      \State Write $n=8k+7$ with $k\ge 2$ and apply the two square-gadget extractions in \Cref{lem:8k+7-splittable}; this recovers both auxiliary polynomials and isolates $P_{2k+1}$.
      \State Recurse on $P_{2k+1}$, thereby reconstructing the internal powers $H_2,H_4$.
      \State Only now decode the recovered auxiliary gadgets by the $\mathcal Q_d$ dispatch of \Cref{lem:odd-gadgets-H2H4} (i.e.\ \Cref{lem:Q4k+1-from-H2}, \Cref{lem:Q-odd-degree-with-powers}, or \Cref{lem:barQ15,lem:barQ8k+7}), with \Cref{alg:decode-Q-2kminus1} for nested Mersenne blocks.
    \EndIf
  \end{algorithmic}
\end{algorithm}

\begin{theorem}[Cost-optimal compatible pairs in odd degree]
\label{thm:odd-realizable-pairs}
    Let $n\ge3$ be odd with $n\ne7$, and assume that $\mathbb F$ is
    $n$-admissible.  Then there exists a splittable pair
    $(T^{(1)}_n,T^{(2)}_n)$ that is compatible on an explicit window with no
    auxiliary data and has a joint $(n-1)/2$-realization
    (\Cref{def:joint-realization}) of multiplicative height at most
    $2\lceil\log_2n\rceil+3$.  The realization
    records a monic quadratic $H_2$ and, when $n\ge5$, a monic quartic $H_4$.
    In particular, $P_n=xT^{(1)}_n+T^{(2)}_n$ is decodable.
\end{theorem}
\begin{proof}
    We argue by strong induction on odd $n$, proving simultaneously the
    asserted compatible-pair conclusion, the exact joint multiplication
    count, and the byproduct statement.  For the height conjunct, the ledger
    of \Cref{thm:construction-height-peeled} threaded through the same
    induction gives $2\lceil\log_2n\rceil+3$ along the pair branches (the
    bound $B(L)=2L+3$ there); the same constant is the height conjunct of
    the machine-checked Lean statement (\texttt{odd\_realizable\_pairs}) for
    the very circuit exhibited here (see the closing remark of
    \Cref{sec:peeled-Q}).
    Since $n$-admissibility implies $m$-admissibility for every $m<n$, all
    recursive pairs and auxiliary gadgets below satisfy their stated
    characteristic hypotheses over the same field.
    The strengthened induction invariant is that every constructed pair records
    its monic quadratic $H_2$ as a polynomial byproduct, and every constructed
    pair of degree at least five also records a monic quartic $H_4$.  These
    polynomials use no extra gates: they are intermediate values already
    displayed in the construction.  The degree-$3$ base records its $H_2$;
    the $4k+1$ and special constructions display both powers; the $8k+3$ step
    inherits $H_2$ and obtains $H_4$ as the byproduct of its
    $Q_{4k+1}(x,H_2)$ call; and the $8k+7$ step inherits both powers from its
    smaller pair.  Thus the byproduct invariant is preserved by every branch.

    Unconditional compatibility is stronger than decodability and is precisely what permits a
    smaller pair to be used before any powers internal to its construction have
    been reconstructed.  The direct proofs in
    \Cref{lem:4k+1-splittable,lem:8k+3-splittable,lem:8k+7-splittable}
    verify the cutoff formulas for each of the three recursive branches; no
    general composition assertion is used.

    The realizability assertion is carried by those same branches, not by a
    separate numerical recurrence.  In either odd induction step, execute the
    smaller joint circuit once, reuse its recorded $H_2,H_4$ wires as inputs
    to the auxiliary gadgets, and then form the two displayed
    difference-of-squares outputs.  Each difference of squares uses the one
    product $(S_3+S_2)(S_3-S_2)$ already shown in the construction algorithm.
    Thus the two output components share the smaller circuit and the power
    byproducts literally; the sums below count the gates of this one joint
    circuit.  The $4k+1$ and finite branches are already displayed as joint
    circuits in the same way.

    The finite cost bases have two distinct causes.  For $n=27$ and
    $n=31$, the relevant recursion would call degree~$7$.  Although the
    septic is algebraically splittable by \Cref{rem:canonical-split}, we do
    not have the three-product compatible joint realization, with recorded
    powers, required by either cost-optimal induction step.  For $n=15$, the smaller degree-$3$
    pair is compatible, but records only $H_2$, whereas the degree-$7$
    auxiliary gadget in the generic $8k+7$ step also requires $H_4$.
    Constructing that quartic separately would exceed the target count; the
    special degree-$15$ circuit instead replaces the generic $Q_3$ block by a
    scalar shift of $H_2$ and spends the saved product on $H_4$.  Thus these
    cases are not artifacts of the decoder proof.

    For $n\equiv 1\pmod 4$ and $n\ge 5$, write $n=4k+1$ and use
    \Cref{lem:4k+1-splittable}.  Its joint-realization clause gives exactly
    $2k=(n-1)/2$ products and records both powers.

    For $n\equiv 3\pmod 8$ we can write $n=8k+3$ with $k\ge 0$.
    If $k=0$ then $n=3$ is supplied by the compatible base
    \Cref{lem:base-three-compatible}.  Hence assume $k\ge1$.
    If $k=3$ then $n=27$ is covered, with its exact cost and recorded powers,
    by \Cref{lem:special-cases-splittable}.
    Otherwise $k\ne 3$, so $2k+1<n$ is odd and $2k+1\ne 7$; by the induction
    hypothesis we have a compatible joint realization for $2k+1$ costing
    $k$ products.  The full conclusion of \Cref{lem:8k+3-splittable} gives
    the required decoder, compatibility, recorded powers, and joint cost
    \[
       k+3k+1=4k+1=\frac{(8k+3)-1}{2}.
    \]

    For $n\equiv 7\pmod 8$ we can write $n=8k+7$ with $k\ge 0$.
    The case $k=0$ is exactly $n=7$ and is excluded.
    If $k=1$ then $n=15$, and if $k=3$ then $n=31$; both are covered, with
    their exact costs and recorded powers, by
    \Cref{lem:special-cases-splittable}.  Otherwise $k\ne1,3$, so $k\ge2$
    and $2k+1<n$ is odd with $2k+1\ne7$.  The induction hypothesis gives a
    compatible joint realization for $2k+1$ costing $k$ products, and
    the full conclusion of \Cref{lem:8k+7-splittable} gives the required
    decoder, compatibility, recorded powers, and joint cost
    \[
       k+3k+3=4k+3=\frac{(8k+7)-1}{2}.
    \]
    The byproduct discussion above verifies in each branch that these are
    realizations of the full strengthened invariant, completing the induction.
\end{proof}

\begin{theorem}[Exact multiplication count]\label{thm:construction-count}
    For every odd $n\ge3$ with $n\ne7$, the construction computes the
    splittable pair $(T^{(1)}_n,T^{(2)}_n)$ using exactly $(n-1)/2$
    multiplications, and therefore computes
    $P_n=xT^{(1)}_n+T^{(2)}_n$ using $(n+1)/2$ multiplications.  The direct
    septic construction uses four multiplications as well.  Consequently, over
    every $n$-admissible field, each monic degree-$n$ polynomial with $n\ge3$
    is covered by a decodable circuit using
    $\lfloor n/2\rfloor+1$ multiplications.
\end{theorem}
\begin{proof}
    For odd $n\ne7$, \Cref{thm:odd-realizable-pairs} supplies the joint
    $(n-1)/2$-realization and its decoder.  Multiplying $T^{(1)}_n$ by $x$
    once gives $P_n$ and raises the cost to $(n+1)/2$.  For $n=7$ the direct
    circuit in \Cref{lem:septic-base} already has the same cost, namely four.

    Finally, if $n\ge4$ is even, first use the odd construction for a monic
    degree-$(n-1)$ polynomial $Q$, at cost $n/2$, and output
    $c_0+xQ$, using one additional product.  This gives $n/2+1$ products.
    Its decoder first reads $c_0$ from the constant coefficient and then
    shifts the remaining coefficients down by one before invoking the decoder
    for $Q$, so this even lift is decodable.  Degrees $1$ and $2$ use
    $x+\alpha_0$ and $(x+\alpha_1)x+\alpha_0$, respectively, and are
    decodable with zero and one product.  Finally,
    \Cref{lem:polynomial-left-inverse-automorphism} turns each of these
    polynomial decoders into a two-sided polynomial preprocessing map.
    Hence the families cover every monic coefficient vector, as asserted.
\end{proof}

\begin{remark}[Exact characteristic restriction]
\label{rem:bad-primes}
    The hypothesis of $n$-admissibility in
    \Cref{thm:odd-realizable-pairs,thm:construction-count} is used only to
    invert the integers listed in \eqref{eq:bad-primes} along the routing
    of \Cref{alg:final-construction}.  Consequently both theorems hold
    verbatim, with the same circuits, the same decoder
    \Cref{alg:final-decoder}, the same multiplication count
    $\lfloor n/2\rfloor+1$ and the same height, over every field whose
    characteristic is not in $\mathrm{Bad}(n)$; equivalently, over every
    field in which each element of $\mathrm{Bad}(n)$ is a unit.  The set
    $\mathrm{Bad}(n)$ is computed by \texttt{tools/bad\_primes.py}, always
    contains~$2$ for $n\ge5$, and its odd elements are at most $(n-1)/4$,
    so characteristic zero or $p>n$ remains a uniform sufficient condition.
    For example $\mathrm{Bad}(n)=\{2\}$ for $5\le n\le12$ and for
    $n=15,17,19,23,31,33,35,39,47$, while $\mathrm{Bad}(13)=\{2,3\}$,
    $\mathrm{Bad}(21)=\{2,5\}$ and $\mathrm{Bad}(29)=\{2,3,7\}$: the
    degree-$47$ circuit is decodable over $\mathrm{GF}(3)$, whereas the
    degree-$13$ crown divides by $2k=6$.  The Lean statement
    \texttt{odd\_realizable\_pairs} certifies the uniform hypothesis that
    $1,\dots,n$ are units (\Cref{sec:formalization-map:units}); the sharper
    set $\mathrm{Bad}(n)$ is read off the displayed pivots and is not
    machine-checked.
\end{remark}

